\section{An explicit grid partition of the chain}
\label{sec:huang_chain_grid}

The following construction supplies the two support layouts in
\cite[Figure 2]{Osborne2006Dynamics}. It fixes the blocks, internal boundaries
and boundary buffers directly from the chain length and two integer parameters.

\begin{definition}[Grid blocks and boundary buffers]
    \label{def:huang_chain_grid}
    \lean{TNLean.MPS.AreaLaw.chainGridBlock,TNLean.MPS.AreaLaw.chainGridCut,
        TNLean.MPS.AreaLaw.chainCutPair,TNLean.MPS.AreaLaw.chainCutBuffer,
        TNLean.MPS.AreaLaw.chainGridOwner}
    \leanok
    Let $L,w,r$ be nonnegative integers, with $w>0$. On the chain
    $\{0,\ldots,L-1\}$, set
    \begin{align*}
        C_j & =[jw,(j+1)w),       & 0\le j & <\lfloor L/w\rfloor+1,  \\
        c_k & =(k+1)w,            & 0\le k & <\lfloor(L-1)/w\rfloor, \\
        S_k & =[c_k-1,c_k+1),     &
        K_k & =[c_k-1-r,c_k+r+1).
    \end{align*}
    All intervals are intersected with the chain; a negative left endpoint is
    replaced by zero. The subtraction $L-1$ in the label count is truncated
    at zero, so for $L=0$ the cut label set is empty.
    A site $x$ is assigned the block label $\lfloor x/w\rfloor$.
\end{definition}

\begin{lemma}[Disjoint blocks covering the chain]
    \label{thm:huang_chain_grid_blocks}
    \lean{TNLean.MPS.AreaLaw.card_chainGridBlock_le,
        TNLean.MPS.AreaLaw.chainGridBlock_width,
        TNLean.MPS.AreaLaw.pairwise_disjoint_chainGridBlock,
        TNLean.MPS.AreaLaw.mem_chainGridBlock_owner,
        TNLean.MPS.AreaLaw.iUnion_chainGridBlock_eq_univ,
        TNLean.MPS.AreaLaw.chainGrid_label_counts,
        TNLean.MPS.AreaLaw.chainGridCut_bounds}
    \leanok
    \uses{def:huang_chain_grid}
    The blocks $C_j$ are pairwise disjoint, each has at most $w$ sites, and
    their union is the chain. Their nominal endpoint widths equal $w$.
    Every site belongs to its assigned block. There are
    $\lfloor L/w\rfloor+1\le L+1$ block labels and
    $\lfloor(L-1)/w\rfloor\le L$ cut labels.
    Each enumerated cut satisfies $0<c_k<L$.
\end{lemma}
\begin{proof}
    \leanok
    The division algorithm gives
    $\lfloor x/w\rfloor w\le x<(\lfloor x/w\rfloor+1)w$.
    Distinct grid intervals do not intersect. The internal-cut label bound
    gives $(k+1)w\le L-1$, and the remaining counts follow from integer division.
\end{proof}

\begin{lemma}[Exact neighborhoods and disjoint boundary buffers]
    \label{thm:huang_chain_grid_buffers}
    \lean{TNLean.MPS.AreaLaw.chainCutBuffer_width_le,
        TNLean.MPS.AreaLaw.card_chainCutBuffer_le,
        TNLean.MPS.AreaLaw.chainNeighborhood_chainCutPair_eq,
        TNLean.MPS.AreaLaw.pairwise_disjoint_chainGridBuffer,
        TNLean.MPS.AreaLaw.chainCutPair_eq_pair}
    \leanok
    \uses{def:huang_chain_grid, thm:huang_chain_grid_blocks}
    For every internal cut, $S_k=\{c_k-1,c_k\}$ and its closed radius-$r$
    neighborhood is exactly $K_k$. The nominal width and cardinality of $K_k$
    are at most $2r+2$. If $2r+2\le w$, the buffers are pairwise disjoint.
\end{lemma}
\begin{proof}
    \leanok
    Taking the closed integer-radius neighborhood expands the adjacent pair
    by $r$ sites on either side, with clipping at the chain ends. Distinct
    grid cuts are separated by at least $w$. When $2r+2\le w$, the right
    endpoint of the earlier buffer is no larger than the left endpoint of
    the later one.
\end{proof}

\begin{lemma}[Classification of nearest-neighbor edges]
    \label{thm:huang_chain_grid_edges}
    \lean{TNLean.MPS.AreaLaw.chainGrid_edge_classification,
        TNLean.MPS.AreaLaw.chainGrid_boundary_edge_eq_pair,
        TNLean.MPS.AreaLaw.chainGrid_pair_not_subset_block,
        TNLean.MPS.AreaLaw.chainGrid_edge_in_block_iff}
    \leanok
    \uses{def:huang_chain_grid, thm:huang_chain_grid_blocks}
    For every edge $\{i,i+1\}$ of the chain, exactly one of the following holds:
    the edge is contained in a grid block, or $i+1=c_k$ for an enumerated
    internal cut. In the second case, the edge is exactly $S_k$ and is
    contained in no single block.
\end{lemma}
\begin{proof}
    \leanok
    Put $j=\lfloor i/w\rfloor$. If $i+1<(j+1)w$, both endpoints lie in $C_j$.
    Otherwise $i<(j+1)w\le i+1$, so $(j+1)w=i+1$ and the cut is internal.
    The pair on either side of a grid boundary cannot lie in one grid interval.
\end{proof}
